\documentclass[11pt,a4paper]{article}
\usepackage[T1]{fontenc}
\usepackage[utf8]{inputenc}
\DeclareUnicodeCharacter{220E}{\ensuremath{\square}}
\usepackage{lmodern,amsmath,amssymb,amsthm,mathtools}
\usepackage[margin=25mm,headheight=14pt]{geometry}
\usepackage{microtype,longtable,booktabs,array,calc,graphicx,enumitem,fancyhdr,xurl}
\usepackage[hidelinks]{hyperref}
\DeclareMathOperator{\Log}{Log}
\DeclareMathOperator{\supp}{supp}
\providecommand{\tightlist}{\setlength{\itemsep}{0pt}\setlength{\parskip}{0pt}}
\newcommand{\passthrough}[1]{#1}
\newtheorem{theorem}{Theorem}
\newtheorem{lemma}[theorem]{Lemma}
\setcounter{secnumdepth}{2}
\setlength{\parindent}{1.2em}
\setlength{\parskip}{3pt plus 1pt}
\setlength{\emergencystretch}{3em}
\setlist{nosep,leftmargin=1.5em}
\pagestyle{fancy}
\fancyhf{}
\fancyhead[L]{\small Correlated translated probes}
\fancyhead[R]{\small D. J. Raymond}
\fancyfoot[C]{\small\thepage}
\renewcommand{\headrulewidth}{0.3pt}
\hypersetup{pdftitle={Correlated translated probes and the bound Lambda < 0.149396 for the de Bruijn-Newman constant},pdfauthor={Dale James Raymond},pdfsubject={Computer-assisted upper bound; version 1.0, pre-deposit editorial revision R2},pdfkeywords={de Bruijn-Newman constant, Riemann xi function, heat flow, logarithmic derivatives, interval arithmetic, zero dynamics}}
\title{Correlated translated probes and the bound\\$\Lambda<0.149396$ for the de Bruijn--Newman constant}
\author{Dale James Raymond\\[3pt]
\small ORCID: \href{https://orcid.org/0009-0001-1682-8529}{0009-0001-1682-8529}\\
\small \href{mailto:dale_raymond@icloud.com}{\texttt{dale\_raymond@icloud.com}}}
\date{21 September 2026 --- Version 1.0}
\begin{document}
\maketitle
\begin{abstract}
A computer-assisted argument gives the upper bound
$\Lambda\le0.149395073827688<0.149396$
for the de Bruijn--Newman constant in the Polymath normalisation.
The construction starts from the classical width-one zero strip at time zero.
A finite-head localisation wall and uniform logarithmic-derivative estimates
are coupled to a decreasing zero-height barrier.
The main refinement retains nonnegative surplus terms in a one-point
value--derivative inequality: their phases at two prescribed translations
cannot simultaneously realise the separate worst cases.
A second gain comes from the finite Euler logarithmic derivative.
The proof accommodates unequal coefficients and denominators, source-frequency
seams, physical cutoff changes and the complete approximation remainder.
The source certificates cover 608 time cells, with six frequency ranges per
cell including an infinite final range. The terminal comparison is enclosed
by exact rational arithmetic. The deposited source and evidence allow the
numerical inequalities to be regenerated. The analytic arguments, including
the larger-height approximation, are provided in a companion text; they are
not formalised by the numerical replay.
\end{abstract}
\noindent\textbf{2020 Mathematics Subject Classification.} Primary 11M26; secondary 65G30.\\
\textbf{Keywords.} de Bruijn--Newman constant; Riemann xi function; heat flow;
logarithmic derivatives; Poisson kernels; interval arithmetic.

\section{Introduction}\label{sec:introduction}
The de Bruijn--Newman constant $\Lambda$ is the threshold at which every zero
of the heat-deformed Riemann xi function is real. In the normalisation used
here, the Riemann hypothesis is equivalent to $\Lambda\le0$; the theorem
of Rodgers and Tao gives $\Lambda\ge0$ \cite{polymath,rodgerstao}.
The present paper concerns an explicit positive upper bound, not the
zero endpoint.

\paragraph{Previous and concurrent bounds.}
Polymath proved $\Lambda\le0.22$ \cite{polymath}.
Platt and Trudgian's finite verification of RH subsequently gave
$\Lambda\le0.2$ \cite[Corollary 2]{pt}.
Among the subsequent public preprints, Gomila announced
$\Lambda\le893927/5000000=0.1787854$ in August 2026
\cite{gomila}, and the author's earlier Zenodo paper reported
$\Lambda<0.163297$ in September 2026 \cite{raymond}.
Gordon's concurrent September 2026 manuscript reports
$\Lambda\le B_G=0.157656619095490606768\ldots$ \cite{gordon}, where
\[
 B_G=\frac{3885632262767861213460393068710302759}
 {24646172707879668706230182733520000000}.
\]
Gordon also starts from the classical strip at time zero, but uses different
density and jet estimates and a partial Lean formalisation whose deduction
retains explicit analytic hypotheses. The recent preprints are cited for
context, not as proof dependencies. No theorem, estimate or numerical
certificate from Gomila's or Gordon's work is used to obtain the bound below;
the required new estimates and computations are supplied here. This statement
concerns the mathematical dependency of the proof, not a historical priority
claim. The comparison sources and their status were checked on
21 September 2026; no exhaustive survey of priority is asserted.

\begin{theorem}[Main result]\label{thm:main}
For the de Bruijn--Newman constant in the normalisation of
Section~\ref{the-source-model-and-its-existing-inequalities},
\[
\Lambda\le
\frac{18674384228461}{125000000000000}
=0.149395073827688<0.149396.
\]
\end{theorem}

The earlier paper \cite{raymond} uses a fixed positive-time row and a
final-time covering. Neither its narrow-strip theorem nor its atlas is
assumed here. Instead, the argument combines a time-zero localisation wall,
all-frequency one-point bounds, a translated-phase surplus, and a continuous
zero-height comparison. The contribution is the quantitative two-probe
improvement with its complete physical and arithmetic accounting, rather
than a new proof of the classical heat-flow framework.

\subsection{Dependencies and verification scope}
The companion \emph{Analytic Foundations} is part of the mathematical
submission. It reproduces the supporting source derivations at their exact
scope. The following map distinguishes published inputs, new analytic
arguments and finite numerical certificates.
\begin{center}
\small
\begin{tabular}{@{}p{0.24\linewidth}p{0.43\linewidth}p{0.25\linewidth}@{}}
\toprule
Ingredient & Precise source & Role \\
\midrule
Initial strip; dynamics; collisions; confinement &
\cite[Proposition 3.1; Theorems 3.2 and 1.5(i)]{polymath} & Published analytic inputs \\
Finite RH through $3000175332800$ &
\cite[Theorem 1]{pt}; Foundations A, Section 2 & Published computational theorem \\
Localisation wall & Foundations A, Section 3; regenerated prefix inequalities & New finite certificate \\
Larger-height approximation and derivative errors & Foundations A, Sections 4--5, derived from \cite[Section 6]{polymath} & Written extension and interval bounds \\
One-point bounds and sharp geometry & Foundations A, Sections 6--12; Foundations B, Sections 1--5 & Supporting analytic and numerical results \\
Correlated gain and final endpoint & Sections~\ref{exact-decomposition-of-the-one-point-cone-surplus}--\ref{coupling-the-average-to-the-global-zero-barrier} and \ref{final-publication-refinement-and-opportunity-audit}; deposited certificates & Present refinement \\
\bottomrule
\end{tabular}
\end{center}

The larger-height approximation is not an unqualified invocation of the
height-restricted numerical statement of \cite[Theorem 1.3]{polymath}.
Its derivation and complete remainders are in Foundations A, Sections 4--5.
Both transcendental calculations use Arb/FLINT \cite{arb}; their different
precision and construction choices do not make them independent libraries.
An additional standard-library checker validates the rational comparison
histories. Numerical regeneration checks the finite inequalities, whereas
the analytic reductions require mathematical review. This preprint has not
undergone external peer review or complete proof-assistant formalisation.

\subsection{Organisation and reproducibility}
Sections~\ref{the-source-model-and-its-existing-inequalities}--\ref{uniform-lower-coefficient-bounds-and-complete-spatial-coverage}
construct the joint source estimate. Sections~\ref{variable-displacement-and-the-retained-kernel-laws}--\ref{coupling-the-average-to-the-global-zero-barrier}
apply it to zero motion. Sections~\ref{computation-reproducibility-and-trust-boundary}--\ref{final-publication-refinement-and-opportunity-audit}
state the implementation, limitations and exact terminal comparison.
\texttt{Proof\_and\_Replay.zip} contains the tested sources, numerical runtime,
certificates and default regeneration command. The separate Verification
Guide distinguishes file-integrity checking from numerical regeneration
and analytic review. The old atlas volumes are not required.

\section{The source model and the one-point estimate}\label{the-source-model-and-its-existing-inequalities}

Use the same normalisation
\[
H_t(z)=\int_0^\infty e^{tu^2}\Phi(u)\cos(zu)\,du,
\qquad H_0(z)=\xi((1+iz)/2)/8.
\]
Here
\[
\Phi(u)=\sum_{n=1}^{\infty}
(2\pi^2n^4e^{9u}-3\pi n^2e^{5u})e^{-\pi n^2e^{4u}},
\qquad
\xi(s)=\tfrac12s(s-1)\pi^{-s/2}\Gamma(s/2)\zeta(s).
\]
The normalisation, the identity $\partial_tH_t=-\partial_z^2H_t$,
and the classical zero dynamics follow the conventions of
\cite{polymath}. All zero counts include multiplicity. The detailed
source maps, normaliser $B_t$, and branches of the logarithm are given in
Analytic Foundations, Part A, Section 4.

Write \(S(x,h,t)=-\Im(H'_t/H_t)(x+ih)\). The reference source has
\(X=6000000185827\), \(N_0=690988\), \(L=\log(x/(4\pi))\), and
\[
 s_U^0=(1+h)/2+tL/4-i(x/2+\pi t/8),\qquad
 s_V^0=(1-h)/2+tL/4+i(x/2+\pi t/8).
\]
Its leading gamma is \(\gamma_0=\gamma(x,t,0)e^{-hL/2}\).
The finite Euler multiplier is
\[
 E_t(s)=\prod_{p\in\{2,3,5\}}(1-b_t(p)p^{-s}),
 \qquad b_t(n)=e^{t\log^2n/4}.
\]
The finite Euler multiplier in this paper is distinct from the connected
exponential normalisers used in the earlier fixed-time programme.

Write
\[
P_N(s)=\sum_{n\le N}b_t(n)n^{-s},
\qquad
N=N(x,t)=\left\lfloor\sqrt{\frac{x}{4\pi}+\frac{t}{16}}\right\rfloor.
\]
All logarithms of positive real numbers are natural logarithms. At each
actual source, the complete leading core uses this actual cutoff:

\[
 Q=E_t(s_U^0)\big(P_N(s_U^0)+\gamma_0P_N(s_V^0)\big)
   =\sum_{j\ge0}a_jz_j,
\]
with \(a_0=1,z_0=1\), all \(a_j\ge0\), and \(|z_j|=1\).
The product head through \(N_0\) is combined first; the full product tail,
reflected branch and every later activated term remain explicit signed
labels (their signs are absorbed in their phases). At fixed \(x,t,N\),
\[
 Q_h=-V/2,\qquad V=\sum_j\lambda_j a_jz_j,
 \quad \lambda_0=0,\quad\lambda_j\ge0.
\]
For first-head index \(n\), \(\lambda_n=\log n\) and
\[
 z_n=\exp\big(i(x/2+\pi t/8)\log n\big).
\]
For reflected label \((d,n)\), \(\lambda=L-\log n+\log d\).

The one-point calculation in Analytic Foundations, Part A, certifies, on each complete time--frequency rectangle,
\[
0\le A_-\le A:=\sum_{j>0}a_j\le A_+<4/5,
\qquad B:=\sum_j\lambda_ja_j\le B_+,
\]
with a rational \(C\in(0,3.4)\) for which
\[
 (1+A_-)[C(1+A_-)-B_+]-D_+>0,
\]
where \(D_+\) bounds the complete negative-pair budget defined below. It also
provides the Euler logarithmic-derivative cap \(E_\ell\), the complete value
and height-derivative remainder caps \(R_0,R_1\), and a strict physical floor
\(\ell\). These are reconstructed, not read solely from a PASS marker.

The physical transfer in Analytic Foundations, Part A, Sections 5 and 9, yields
\[
 S>\ell+
 \frac{C|Q|^2-\Re(V\overline Q)}{2|Q|^2}
 +\left(E_\ell-\Re\frac{(E_t(s_U^0))_h}{E_t(s_U^0)}\right).
 \tag{1}
\]
Here \(\ell\) is rounded down from
\[
 L_{\rm lower}/4-10^{-6}-E_\ell-C/2
 -\frac{R_1+B_+R_0/[2(1-A_+)]}{1-A_+-R_0}.
\]
Equation (1) is the one-point argument with the discarded nonnegative terms
retained. Both complete physical remainders are still charged. Their caps do
not depend on the discarded phases. The leading replacement and its height
derivative are justified before using (1); no unknown error is differentiated.

\section{The nonnegative surplus in the one-point estimate}\label{exact-decomposition-of-the-one-point-cone-surplus}

Set
\[
 w_{ij}=a_i a_j(\lambda_i+\lambda_j-2C),\qquad
 D_C=2\sum_{i<j,\ w_{ij}<0}|w_{ij}|.
\]
\begin{lemma}[Surplus decomposition]\label{lem:surplus}
With the notation above, the following identity holds. Each term in its
last two sums is nonnegative.

\[
\begin{split}
 H_C&:=C|Q|^2-\Re(V\overline Q)\\
 &=(1+A)[C(1+A)-B]-D_C\\
 &\quad+\sum_{w_{ij}\ge0}w_{ij}[1-\Re(z_i\overline z_j)]\\
 &\quad+\sum_{w_{ij}<0}|w_{ij}|[1+\Re(z_i\overline z_j)].
\end{split}\tag{2}
\]
\end{lemma}
\begin{proof}
Expand $|Q|^2$ and $\Re(V\overline Q)$, group the terms with indices $i<j$,
and separate the signs of $w_{ij}$. For $w_{ij}<0$, subtracting
$2|w_{ij}|$ gives the displayed form. Nonnegativity follows from
$-1\le\Re(z_i\overline z_j)\le1$.
\end{proof}
The one-point argument proves the first
line after subtraction of the \emph{complete} \(D_C\) is positive. It follows that
any selected subset of the displayed surplus terms is a valid lower bound
for \(H_C\).

The implementation selects only constant--head pairs \((0,n)\),
\(2\le n\le1024\). All other surplus terms are left at their proven lower
bound zero. In particular, no omitted negative coefficient is discarded:
its entire worst-case debit remains in the one-point cone certificate.

The denominators at the two probes need not be equal. Pointwise,
\[
 \frac{H_C}{2|Q|^2}\ge
 \frac{\text{selected nonnegative surplus}}{2(1+A_+)^2}.
 \tag{3}
\]
The direction in (3) uses nonnegativity. Averaging the two quotients first and
then replacing them by a quotient of averages would in general be false.

\section{The translated-phase inequality}\label{why-two-translated-sources-cannot-both-saturate-these-terms}

For any unit complex numbers \(z,u\),
\[
 \frac{\Re(zu)+\Re(z\overline u)}2=\Re z\,\Re u.
\]
\begin{lemma}[Paired-phase lower bound]\label{lem:paired}
For $|z|=|u|=1$, $\epsilon\in\{-1,1\}$ and
$c_+,c_-\ge c_0\ge0$, the following inequality holds.
\[
\frac{c_+[1-\epsilon\Re(zu)]
+c_-[1-\epsilon\Re(z\overline u)]}{2}
\ge c_0(1-|\Re u|).
\tag{4}
\]
\end{lemma}
This proves a deterministic gain with arbitrary common centre phase. There
is no probability, generic-phase or averaging-in-frequency assumption.

At the actual two points \(x\pm a\), the first-head phases obey exactly
\[
 z_n^\pm=z_n^c e^{\pm ia\log n/2}.
\tag{5}
\]
This is an identity for the specified leading source. It continues to hold
when the physical truncations at the two points differ; every selected index
is at most 1024, far below both truncations. The nonlinear source replacement
errors have already been included separately at both physical points in (1).
No shared phase is imposed on the gamma or high-tail labels.

Let the two points lie in source frequency rectangles \(j,k\), with certified
cone constants \(C_j,C_k\), amplitude caps \(A_j,A_k\), and strict physical
floors \(\ell_j,\ell_k\). Let \(\underline a_{j,n}\) bound the first-head
amplitude below on rectangle \(j\). Define
\[
 g_{jk,n}=\max\{0,\log n-2\max(C_j,C_k),
                       2\min(C_j,C_k)-\log n\},
\]
\[
 m_{jk,n}=g_{jk,n}\min\left\{
 \frac{\underline a_{j,n}}{2(1+A_j)^2},
 \frac{\underline a_{k,n}}{2(1+A_k)^2}\right\}.
\tag{6}
\]
If \(\log n\) lies between the two thresholds, (6) is zero and no credit is
taken. Otherwise both surplus terms in (2) have the same sign convention;
(3)--(5) give the certified moment-cone contribution
\[
 \boxed{G^{\rm cone}_{jk}(a)
 =\sum_{n=2}^{1024}m_{jk,n}
     \left(1-\left|\cos\frac{a\log n}{2}\right|\right).}
 \tag{7}
\]
This includes recoverable surplus from negative-budget pairs as well as
positive pairs. It does not reduce the verified one-point budget \(D_+\).

\section{A second gain from the finite Euler logarithmic derivative}\label{a-second-gain-from-the-finite-euler-logarithmic-derivative}

Put
\[
 \rho_p=b_t(p)p^{-[(1+h)/2+tL/4]},\qquad
 \vartheta_p=(x/2+\pi t/8)\log p.
\]
Since \(\rho_p<1\), absolutely convergent expansion gives
\[
 \Re\frac{E_h}{E}
 =\frac12\sum_{p=2,3,5}\log p\sum_{r\ge1}
        \rho_p^r\cos(r\vartheta_p).
\]
The complete phase-free cap minus this expression is the nonnegative series
\[
 \frac12\sum_p\log p\sum_{r\ge1}
        \rho_p^r[1-\cos(r\vartheta_p)].
\tag{8}
\]
The certified Euler cap is at least this full phase-free cap at every source.
Applying (4), keeping only the first eight harmonics, gives
\[
 \boxed{G^{\rm Euler}_{jk}(a)=
 \frac12\sum_{p=2,3,5}\log p\sum_{r=1}^8
 \min(\underline\rho_{j,p}^r,\underline\rho_{k,p}^r)
 \left(1-\left|\cos\frac{ra\log p}{2}\right|\right).}
 \tag{9}
\]
All higher harmonics remain covered: their contributions to (8) are
nonnegative. They are not bounded by an unjustified zero remainder.

The cone and Euler gains may be added. They arise from two different explicit
addends in (1); no common coefficient resource is allocated twice inside
one identity. The underlying phases can be dependent--the sum of two valid
lower bounds requires no independence.

\begin{theorem}[Uniform two-probe improvement]\label{thm:joint}
Assume the complete one-point bounds on source rectangles $j,k$ stated in
Section~\ref{the-source-model-and-its-existing-inequalities}. With the
lower coefficients in (6)--(9), the two specified translated sources satisfy

\[
\boxed{
\frac{S(x+a,h,t)+S(x-a,h,t)}2
>
\frac{\ell_j+\ell_k}{2}
+G^{\rm cone}_{jk}(a)+G^{\rm Euler}_{jk}(a).
}\tag{10}
\]
\end{theorem}
\begin{proof}
Apply Lemma~\ref{lem:paired} to the selected terms in
Lemma~\ref{lem:surplus} after the pointwise denominator estimate (3).
Apply the same inequality to each of the retained Euler harmonics in (8),
and add the two lower bounds in (1). This gives (10).
\end{proof}
All physical approximation and derivative errors are unchanged from the
complete one-point inputs. In particular, (10) does not assert that either
individual probe satisfies the improved average floor.

\section{Uniform lower coefficient bounds and complete spatial coverage}\label{uniform-lower-coefficient-bounds-and-complete-spatial-coverage}

The raw Euler-product head coefficient is
\[
 c_n(t)=\sum_{d\mid(n,30)}\mu(d)
        \left(\prod_{p\mid d}b_t(p)\right)b_t(n/d).
\]
For \(n\le N_0\) it is nonnegative and nondecreasing in \(t\ge0\).
Indeed the finite heat graph expansion retains precisely graphs in which
all marked selected-prime occurrences are nonisolated. Every retained
weight and its heat factor is nonnegative and nondecreasing. Repeated prime
occurrences are included. At the three selected primes the coefficient is
identically zero.

For a time cell \([t_L,t_R]\) and finite frequency rectangle with upper bound
\(L_j^+\), valid lower bounds are
\[
 \underline a_{j,n}
 =c_n(t_L)\exp\left[-\left(\frac{1+h}{2}
                   +\frac{t_RL_j^+}{4}\right)\log n\right],
\tag{11}
\]
\[
 \underline\rho_{j,p}
 =\exp\left[-\frac{1+h}{2}\log p
        +\frac{t_R}{4}\log p(\log p-L_j^+)\right].
\tag{12}
\]
Equation (11) does not assume monotonicity of the product of these factors;
it bounds the increasing coefficient and decreasing exponential separately.
For (12), \(L>\log p\), so the exact weight decreases in both \(t\) and \(L\).
On the infinite last rectangle, use lower bound zero. Its full previously
proved source floor, including every future state activation, is retained.

The source frequency variable is \(s=L/2-\log N_0\), with endpoints
\[
0,\;1/20,\;3/20,\;2/5,\;1,\;2,\;\infty.
\]
Both translated arguments satisfy \(x\pm a\ge X-2\). The calculation checks
\(X-4>4\pi N_0^2\), so they are in the declared source range. The complete
distance between them is less than four, whereas every finite frequency
interval has physical width larger than \(10^{11}\). Thus the possible
ordered pairs are exactly covered by
\[
(j,j),\ 0\le j\le5,\qquad (j,j+1),\ 0\le j\le4.
\tag{13}
\]
The lower bound used in zero motion is the minimum of (10) over all eleven
possibilities. This includes points arbitrarily close to a frequency seam,
not just centres away from seams. Physical truncation changes require no
additional correlation assumption: the selected low labels occur on both
sides, while all other terms and complete state-pair errors remain in the
individual source certificates.

\section{Variable displacement and the retained kernel laws}\label{variable-displacement-and-the-retained-kernel-laws}

For each catalogue displacement \(a=rh\), the new gains are computed at that
exact rational displacement. There are ten distinct ratios in the retained
thirteen-kernel catalogue.

For the sharp continuum law
\[
 a(y,h)=\sqrt{(h^2-y^2)/3},\qquad
 v=y^2/h^2,\qquad 29v^2-v-1\le0,
\]
the new supplier uses one conservative displacement interval
\[
 h\sqrt{(1-51/250)/3}\le a\le h/\sqrt3.
\tag{14}
\]
The exact positive root \(v_*\) is less than \(51/250\). Thus (14) contains
every displacement at which the sharp kernel can be selected. Each factor
\(1-|\cos(a\lambda/2)|\) is replaced by its interval lower bound on (14).
The bound includes interior extrema of cosine; it is not interpolation
between endpoint cosines. The new calculation is deliberately not optimised
over a smaller attained displacement interval.

All displacements are less than two, within the spatial padding in Analytic Foundations, Part A. They are
held fixed when differentiating the probe in height. Their possible
functional dependence on the hypothetical contact height creates no extra
derivative term: the zero-motion comparison evaluates the kernel inequality
instantaneously and does not differentiate a moving probe path.

\section{Coupling the average to the global zero barrier}\label{coupling-the-average-to-the-global-zero-barrier}

For a highest simple upper zero of height \(y\), write \(q=y^2\). The retained
kernel comparison and the complete own-pair subtraction give
\[
 q'\le-2-\frac{2\lambda q}{h}
       \left[\frac{S(x+a,h,t)+S(x-a,h,t)}2
                         -S_{\rm self}(q)\right].
\tag{15}
\]
Its exact self-pair contribution is
\[
 S_{\rm self}(q)=
 \frac{2h(A+h^2-q)}{(A+h^2-q)^2+4Aq},\qquad A=a^2.
\]
For the sharp continuum kernel, \(\lambda=2\) and
\[
 S_{\rm self}(q)=\frac6{h(4-q/h^2)}.
\]
No self-pair credit is taken twice. A stronger bound for the average in
(10) directly strengthens the coefficient in (15).

On every comparison step, the exact consumer chooses only a kernel whose
height-ratio restriction is satisfied by the upward enclosure of \(q\).
The source floor is the matching \emph{joint} floor for that displacement, not
an individually improved probe floor. The self-pair upper bound is checked
on the entire permitted decreasing range of \(q\). Thus a constant rate
\(c>0\) gives \(q'<-2-cq\) through the step.

The source schedule has 608 cells of width \(1/4000\), with unchanged rational
probe heights. The archived joint-probe baseline used four comparison steps
per cell. The publication refinement uses 128 steps per cell, as proved and
checked in Section~\ref{final-publication-refinement-and-opportunity-audit}. The initial bound is \(q_0=(1.000001)^2>1\); the final positive height
is \(q_e=10^{-8}\), followed by de Bruijn contraction costing \(q_e/2\).
Upward dyadic enclosures of each exponential solution are used. A distinct
rational exponential/logarithm implementation supplies a second endpoint.

The global first-contact argument is given in Analytic Foundations,
Part A, Section 11, using the zero-dynamics and confinement results of
\cite[Proposition 3.1 and Theorem 1.5(i)]{polymath}:
the wall excludes nonreal zeros from the finite head, positive-time
confinement makes first contact attainable, backward Hermite splitting
excludes multiple contact, and the simple-zero inequality contradicts first
contact. At a time knot, the left-hand row is used. The initial strict margin
is checked numerically, so compactness is used only after positive time.
The proof does not infer a smaller-time strip by reversing heat flow.

\section{Numerical verification and reproducibility}\label{computation-reproducibility-and-trust-boundary}

The new arithmetic lower bounds use two constructions of every raw head
coefficient through 1024: complete Euler-divisor inclusion--exclusion, and
descending multiplication by the three finite Euler factors. All subsequent
transcendental values and cosine extrema are enclosed with Arb/FLINT\cite{arb} at
192 and 256 bits. Coefficients, denominator reciprocals and attenuation
factors are rounded down on the \(2^{-80}\) grid. Their sums are then
calculated with exact Python integers. Zero may replace a negative lower
endpoint only where the analytic positivity theorem has established it.

The two lanes produce the same finite-grid hashes. Final acceptance retains
the weaker of the two joint lower bounds. Both new lanes use the same
source-row selection from independently validated one-point rows: the stronger
valid complete enclosure can be selected at each matching rectangle. This
is not a claim that the coarser one-point lane independently proves the final
headline. Both full one-point lanes, their remainders and all domain conditions
are regenerated and checked.

All eleven spatial pair classes, all 608 time cells and all eleven translation
modes are required. The last spatial interval is infinite and retains its
full old floor even though its additional lower gains are zero. The exact
checker reconstructs all scalar gain sums, matching domains and minima.
With four steps per cell and the gains disabled, the archived baseline reproduces
the exact preceding endpoint. Its separate cone-only and Euler-only comparisons
measure the effects without pretending their endpoint improvements add linearly.
The final release changes only the scalar subdivision, not these source gains.

Supplementary direct paired checks include an ordinary source, a frequency
seam and a physical truncation seam. They use the actual nonlinear source
coordinates, full two-branch sums, actual cutoffs and complete derivative
error. They are not acceptance premises for the uniform source theorem.

The replay checks integrity and regenerated numerical inequalities. It does
not formally verify the infinite analytic implications or reprove the
published finite-RH computation. Neither of the new lanes is an independent
transcendental library. Exact identities, interval rounding and the scalar
ODE enclosure still depend on their stated arithmetic implementations.

\section{Contribution and limitations}\label{what-is-new-and-what-is-not-claimed}

The new source gain is not obtained from a favourable physical sample or a
conjecture that two worst cases are unlikely. It follows from the exact
relation between selected phases at two specified translations and from
nonnegative surplus terms already present in the complete one-point proof.
The physical denominators may differ. The gamma phase and high branch phases
are not synchronised by assumption. All their previous errors remain.

No product of two million-term polynomials is expanded. Conversely, this is
not the optimal joint two-point estimate: only constant--head pairs through
1024 and eight harmonics of each selected Euler factor are used. Correlation
among other head labels, the reflected branch and the two denominators is
not credited beyond the already proved one-point bounds. Enlarging the
selected positive-surplus set can only help the analytic lower bound, but
whether its certified numerical gain is worthwhile must be evaluated.

The time--height schedule and kernel catalogue are not globally optimised.
The result does not establish \(\Lambda<0.09\) or RH, and no exhaustive
world-record or priority claim is made. It is a further checked corollary
obtained from the analytic estimates and numerical certificates supplied here, with a uniform two-point improvement.

\section{The exact terminal comparison}\label{final-publication-refinement-and-opportunity-audit}

The only final mathematical change to the archived \(\Lambda<0.149399\)
calculation is the number of comparison substeps. The schedule, all source
floors, the selected head indices through 1024, the eight Euler harmonics,
the kernel catalogue, the spatial padding, the self-pair subtraction and
all physical-error contracts are unchanged.

Every source lower bound is valid throughout its full time cell.
At each new substep, let \(q\) be an upward enclosure of the ideal comparison
height. Only a kernel valid on the entire interval from zero to that height
is admissible. Its matching translated-source floor and complete self-pair
cap give a rational \(c>0\) with \(q'<-2-cq\) throughout the substep.
Refining the time partition therefore requires no new functional hypothesis.

For a substep of width \(\Delta t\), the exact solution is
\[
q_{\mathrm{next}}=(q+2/c)e^{-c\Delta t}-2/c.
\]
One construction uses the even alternating Taylor upper bound through degree
18 and 96-bit upward dyadic rounding. The second uses the reciprocal of the
positive exponential Taylor sum through degree 12 and 112-bit upward rounding.
Both enclose the same continuous comparison, not a discretely jumping zero
trajectory.

If the step reaches \(q_e=10^{-8}\), put
\[
z=\frac{c(q-q_e)}{2+cq_e}.
\]
The remaining time is \(\log(1+z)/c\). Each construction bounds this logarithm
upward, checks that the terminal time remains inside the certified source cell,
adds \(q_e/2\), and rounds the final time upward.

The fixed-source audit was:

\begin{longtable}[]{@{}
  >{\raggedleft\arraybackslash}p{(\columnwidth - 4\tabcolsep) * \real{0.3333}}
  >{\raggedleft\arraybackslash}p{(\columnwidth - 4\tabcolsep) * \real{0.3333}}
  >{\raggedleft\arraybackslash}p{(\columnwidth - 4\tabcolsep) * \real{0.3333}}@{}}
\toprule\noalign{}
\begin{minipage}[b]{\linewidth}\raggedleft
Substeps per source cell
\end{minipage} & \begin{minipage}[b]{\linewidth}\raggedleft
Primary exact upper time, displayed
\end{minipage} & \begin{minipage}[b]{\linewidth}\raggedleft
Alternative exact upper time, displayed
\end{minipage} \\
\midrule\noalign{}
\endhead
\bottomrule\noalign{}
\endlastfoot
4 & 0.149398159734722 & 0.149398159735734 \\
8 & 0.149396459831885 & 0.149396459832690 \\
16 & 0.149395675207179 & 0.149395675207190 \\
32 & 0.149395321833941 & 0.149395321833950 \\
64 & 0.149395152324926 & 0.149395152324926 \\
128 & 0.149395073827688 & 0.149395073827688 \\
\end{longtable}

The final two constructions, at 128 substeps per source cell, both give
\[
\boxed{U=\frac{18674384228461}{125000000000000}
=0.149395073827688.}
\]
Compared with the previous exact value \(0.149398159734722\), the decrease is
exactly
\[
0.000003085907034.
\]
These are exact comparisons, not floating-point optimisation scores.

A separate history checker reconstructs every rational rate and time join,
verifies the full quadratic self-pair inequality on the permitted height
interval, and checks each upward endpoint with an even degree-20 exponential
bound. It independently validates terminal logarithms with the odd degree-13
upper series. It rejects modified source floors, rates, time joins, endpoint
ceilings and terminal lags.

\subsection{Scope of the numerical refinement}\label{scope-of-the-opportunity-audit}

This is a bounded refinement of the scalar consumer of already established
functional estimates. It is not an exhaustive optimisation over probe heights,
kernel families, selected surplus terms, or source schedules. Further such
changes may improve the bound, but require their own complete verification.

The release is frozen at the safe public ceiling \(\Lambda<0.149396\).
No optimality or worldwide priority claim follows from this audit.
The number of internal computational iterations is not the publication
version: this separate paper's first prepared public release is version 1.0.

\section{Data availability and author disclosure}\label{availability-and-assistance}

The deposited materials comprise this manuscript, Analytic Foundations,
the numerical programs and runtime, the complete interval enclosures,
exact comparison histories, verification instructions and integrity manifests.
The default regeneration command is given in the Verification Guide.
It recomputes the numerical inequalities; an integrity-only command checks
file identities and does not verify the mathematics. The scientific evidence
archive is unchanged by this editorial revision.

The numerical dependencies are identified by content hashes in the deposit.
The one-point source and sharp-geometry proofs are reproduced in Analytic
Foundations rather than treated as inaccessible external assumptions.
The mathematical inputs from other authors are the published results cited
in the dependency table in Section~\ref{sec:introduction}. Neither
\cite{gomila} nor \cite{gordon} supplies a theorem, estimate or numerical
certificate used to obtain the present bound.

\paragraph{Author disclosure.}
AI assistance was used in mathematical exploration, programming, numerical
experimentation, checking and drafting. The author takes responsibility for
the claims, their presentation and the deposit. Automated review and
numerical replay are not external peer review. No complete proof-assistant
formalisation is claimed. The earlier and concurrent work cited in the
introduction was considered in preparing the comparison and exposition;
independence of proof inputs is not a claim of historical priority or of
absence of awareness of that work.

\paragraph{Publication status.}
This is a preprint prepared for independent mathematical and computational
review. No journal acceptance or exhaustive record-ranking claim is made.
The DOI assigned to this separate record should be used for this manuscript;
the DOI of the earlier paper in \cite{raymond} is not its identifier.
\begin{thebibliography}{9}
\setlength{\itemsep}{5pt}
\bibitem{polymath}
D. H. J. Polymath,
\emph{Effective approximation of heat flow evolution of the Riemann
$\xi$ function, and a new upper bound for the de Bruijn--Newman constant},
Research in the Mathematical Sciences \textbf{6} (2019), article 31.
\href{https://doi.org/10.1007/s40687-019-0193-1}{doi:10.1007/s40687-019-0193-1}.
Numbered references follow
\href{https://arxiv.org/abs/1904.12438v2}{arXiv:1904.12438v2}.

\bibitem{pt}
D. J. Platt and T. S. Trudgian,
\emph{The Riemann hypothesis is true up to $3\cdot10^{12}$},
Bulletin of the London Mathematical Society \textbf{53} (2021), 792--797.
\href{https://doi.org/10.1112/blms.12460}{doi:10.1112/blms.12460};
\href{https://arxiv.org/abs/2004.09765v1}{arXiv:2004.09765v1}.
Theorem 1 supplies finite RH; Corollary 2 gives $\Lambda\le0.2$.

\bibitem{rodgerstao}
B. Rodgers and T. Tao,
\emph{The de Bruijn--Newman constant is non-negative},
Forum of Mathematics, Pi \textbf{8} (2020), e6.
\href{https://doi.org/10.1017/fmp.2020.6}{doi:10.1017/fmp.2020.6};
\href{https://arxiv.org/abs/1801.05914}{arXiv:1801.05914}.

\bibitem{gomila}
J. Gomila,
\emph{A computer-assisted proof of the bound
$\Lambda\le893927/5000000=0.1787854$ for the de Bruijn--Newman constant},
public repository manuscript, August 2026; not peer reviewed at the date
consulted. Repository \texttt{judegomila/dbn-lambda-01787854-candidate-audit},
revision \nolinkurl{a74738deb6d5e0f76887cb36901da08b68dca705},
consulted 21 September 2026.
\href{https://github.com/judegomila/dbn-lambda-01787854-candidate-audit/tree/a74738deb6d5e0f76887cb36901da08b68dca705}{Pinned repository};
\href{https://www.judegomila.com/posts/riemann-lambda-0.1787854}{author's announcement, 19 August 2026}.
Cited for context, not as a proof input.

\bibitem{raymond}
D. J. Raymond,
\emph{Phase-correlated certificates and the bound $\Lambda<0.163297$
for the de Bruijn--Newman constant}, Zenodo preprint, 8 September 2026.
\href{https://doi.org/10.5281/zenodo.22626161}{doi:10.5281/zenodo.22626161}.
Its fixed-time narrow strip and final-time atlas are not premises here.

\bibitem{gordon}
S. Gordon,
\emph{A computer-assisted upper bound for the de Bruijn--Newman constant},
public manuscript and repository, September 2026; not externally peer
reviewed at the date consulted. Repository
\texttt{stefangordon/dbn-upper-bound}, revision
\nolinkurl{64c47f2db67288beb62c24d7a0f7aedd645dc4b3},
consulted 21 September 2026.
\href{https://github.com/stefangordon/dbn-upper-bound/blob/64c47f2db67288beb62c24d7a0f7aedd645dc4b3/manuscript/main.md}{Pinned manuscript};
\href{https://github.com/stefangordon/dbn-upper-bound/blob/64c47f2db67288beb62c24d7a0f7aedd645dc4b3/VERIFICATION.md}{verification scope}.
The partial Lean formalisation retains explicit analytic hypotheses.
Cited for context, not as a proof input.

\bibitem{arb}
F. Johansson,
\emph{Arb: efficient arbitrary-precision midpoint-radius interval arithmetic},
IEEE Transactions on Computers \textbf{66} (2017), 1281--1292.
\href{https://doi.org/10.1109/TC.2017.2690633}{doi:10.1109/TC.2017.2690633};
\href{https://arxiv.org/abs/1611.02831}{arXiv:1611.02831}.

\end{thebibliography}
\end{document}
